Het programma moet controleren of de computer mogelijk een probleem heeft.

Gebruik hiervoor \texttt{if}, \texttt{elif} en \texttt{else}.

Gebruik de volgende waarden:

\begin{table}[!h]
\begin{center}
\caption{CPU Gebruik}
\label{table-cpu}
\begin{tabular}{| c | c |} 
\hline
CPU-gebruik & Status \\
\hline
Lager dan 70\% & OK \\ 
\hline
70\% t/m 90\% & Waarschuwing \\
\hline
Hoger dan 90\% & Kritiek \\
\hline
\end{tabular}
\end{center}
\end{table}

\begin{table}[!h]
\begin{center}
\caption{Geheugengebruik}
\label{table-mem}
\begin{tabular}{| c | c |} 
\hline
Geheugengebruik & Status \\
\hline
Lager dan 70\% & OK \\ 
\hline
70\% t/m 90\% & Waarschuwing \\
\hline
Hoger dan 90\% & Kritiek \\
\hline
\end{tabular}
\end{center}
\end{table}


\begin{table}[!h]
\begin{center}
\caption{Diskgebruik}
\label{table-dsk}
\begin{tabular}{| c | c |} 
\hline
Diskgebruik & Status \\
\hline
Lager dan 80\% & OK \\ 
\hline
80\% t/m 90\% & Waarschuwing \\
\hline
Hoger dan 90\% & Kritiek \\
\hline
\end{tabular}
\end{center}
\end{table}

Let hier op het percentage \textbf{gebruikte} ruimte.

